\documentclass[a4paper,12pt]{article}

\usepackage[T2A]{fontenc}
\usepackage[utf8]{inputenc}
\usepackage[russian]{babel}

\usepackage{amsmath,amssymb,mathtools}
\usepackage{helvet}
\renewcommand{\familydefault}{\sfdefault}
\usepackage{icomma}

\usepackage[a4paper,margin=2.2cm]{geometry}
\usepackage{microtype}
\usepackage{tikz}
\usetikzlibrary{calc}

\usepackage{tabularx,booktabs}
\usepackage{enumitem}
\usepackage{hyperref}
\usepackage{parskip}
\usepackage{fancyhdr}
\usepackage{setspace}
\usepackage{xcolor}

% ------------------------------------------------
% Цветовая палитра
% ------------------------------------------------
\definecolor{accent}{HTML}{008080}
\definecolor{darkaccent}{HTML}{004D4D}
\definecolor{textgray}{HTML}{333333}
\definecolor{softgray}{HTML}{F2F4F4}

% ------------------------------------------------
% Общая типографика
% ------------------------------------------------
\color{textgray}
\onehalfspacing
\setlength{\parindent}{0pt}
\setlength{\parskip}{0.55em}

\hypersetup{
    colorlinks=true,
    linkcolor=darkaccent,
    urlcolor=darkaccent,
    pdftitle={Чек-лист: среднее арифметическое и математическое моделирование},
    pdfauthor={Лёвин Артём Александрович}
}

\setlist[itemize]{
    leftmargin=1.7em,
    itemsep=0.3em,
    topsep=0.3em
}

\setlist[enumerate]{
    leftmargin=2em,
    itemsep=0.7em,
    topsep=0.4em
}

% ------------------------------------------------
% Колонтитулы
% ------------------------------------------------
\pagestyle{fancy}
\fancyhf{}
\fancyhead[L]{\small\color{darkaccent}Среднее арифметическое}
\fancyhead[R]{\small\color{darkaccent}Надя Клименко}
\fancyfoot[C]{\small\thepage}
\renewcommand{\headrulewidth}{0.4pt}
\renewcommand{\headrule}{%
    \hbox to\headwidth{\color{accent}\leaders\hrule height \headrulewidth\hfill}
}

% ------------------------------------------------
% Служебные команды
% ------------------------------------------------
\newcommand{\topicsection}[1]{%
    \vspace{0.7em}
    {\Large\bfseries\color{darkaccent}#1}\par
    \vspace{0.15em}
    {\color{accent}\rule{\linewidth}{0.7pt}}
    \vspace{0.35em}
}

\newcommand{\subtopic}[1]{%
    \vspace{0.4em}
    {\large\bfseries\color{darkaccent}#1}\par
}

\newcommand{\important}[1]{%
    \vspace{0.35em}
    \noindent
    \begin{minipage}{\linewidth}
        \color{darkaccent}\bfseries #1
    \end{minipage}
    \vspace{0.25em}
}

\begin{document}

% ================================================================
% ТИТУЛЬНЫЙ ЛИСТ
% ================================================================
\begin{titlepage}
\thispagestyle{empty}

\begin{center}

\vspace*{2.3cm}

{\color{darkaccent}\rule{0.33\linewidth}{1pt}}

\vspace{1.1cm}

{\Huge\bfseries Чек-лист}

\vspace{0.65cm}

{\LARGE\bfseries
Среднее арифметическое\\[0.2em]
и математическое моделирование}

\vspace{1.1cm}

{\large
Формулы, восстановление суммы,\\
объединение групп и последовательные числа}

\vfill

{\large\bfseries
Лёвин Артём Александрович\\[0.3em]
эксклюзивно для Нади Клименко}

\vspace{0.8cm}

{\large 6 класс}

\vspace{0.5cm}

{\large \today}

\vfill

\begin{tikzpicture}[x=1cm,y=1cm]
    \draw[
        color=accent,
        line width=1.2pt
    ]
    (-6.5,0)
    .. controls (-4.5,0.9) and (-2.3,-0.65) ..
    (0,0)
    .. controls (2.3,0.65) and (4.5,-0.9) ..
    (6.5,0);
\end{tikzpicture}

\vspace{0.7cm}

\end{center}

\end{titlepage}

% ================================================================
% ОСНОВНОЙ ЧЕК-ЛИСТ
% ================================================================

\topicsection{1. Среднее арифметическое: определение}

Среднее арифметическое нескольких чисел получают так:

\[
\boxed{
\overline{x}
=
\frac{x_1+x_2+\dots+x_n}{n}
}
\]

где

\[
x_1,x_2,\dots,x_n
\]

--- данные числа, а \(n\) --- их количество.

\important{Главная идея: сначала складываем все числа, затем делим сумму на их количество.}

\subtopic{Для двух чисел}

\[
\boxed{
\overline{x}=\frac{x_1+x_2}{2}
}
\]

Например, если первое число равно \(7\), а среднее арифметическое двух чисел равно \(5{,}3\), то

\[
\frac{7+x}{2}=5{,}3.
\]

Отсюда

\[
7+x=10{,}6,
\]

\[
\boxed{x=3{,}6}.
\]

\subtopic{Для трёх чисел}

\[
\boxed{
\overline{x}=\frac{x_1+x_2+x_3}{3}
}
\]

\topicsection{2. Восстановление суммы по среднему арифметическому}

Из формулы

\[
\overline{x}
=
\frac{x_1+x_2+\dots+x_n}{n}
\]

можно выразить сумму чисел:

\[
\boxed{
x_1+x_2+\dots+x_n=n\overline{x}
}
\]

Это один из важнейших приёмов занятия.

\important{Если известны среднее арифметическое и количество чисел, всю сумму можно найти одним умножением.}

Например, среднее арифметическое \(14\) чисел равно \(4{,}5\). Тогда их сумма:

\[
S_1=14\cdot4{,}5=63.
\]

Если среднее арифметическое других \(6\) чисел равно \(2{,}75\), то

\[
S_2=6\cdot2{,}75=16{,}5.
\]

\topicsection{3. Среднее арифметическое объединённых групп}

Пусть:

\[
n_1 \text{ чисел имеют среднее } \overline{x}_1,
\]

а

\[
n_2 \text{ чисел имеют среднее } \overline{x}_2.
\]

Сначала восстанавливаем суммы:

\[
S_1=n_1\overline{x}_1,
\qquad
S_2=n_2\overline{x}_2.
\]

После объединения групп общее среднее равно

\[
\boxed{
\overline{x}
=
\frac{n_1\overline{x}_1+n_2\overline{x}_2}
     {n_1+n_2}
}
\]

Для рассмотренного примера:

\[
\overline{x}
=
\frac{14\cdot4{,}5+6\cdot2{,}75}{14+6}
=
\frac{63+16{,}5}{20}.
\]

Следовательно,

\[
\boxed{\overline{x}=3{,}975}.
\]

\important{Среднее двух групп нельзя просто находить как среднее двух средних, если количества чисел в группах различаются.}

Например,

\[
\frac{4{,}5+2{,}75}{2}
\]

в данной задаче использовать нельзя, поскольку в первой группе \(14\) чисел, а во второй --- \(6\).

\topicsection{4. Математическая модель}

В текстовой задаче полезно сначала перевести условие на математический язык.

Последовательность работы:

\[
\boxed{
\text{текст}
\longrightarrow
\text{обозначения}
\longrightarrow
\text{формула}
\longrightarrow
\text{уравнение}
\longrightarrow
\text{ответ}
}
\]

Например:

\[
\text{«среднее двух чисел равно }42\text{»}
\]

переводится в

\[
\frac{x_1+x_2}{2}=42.
\]

Если известно, что второе число в \(2{,}5\) раза больше первого, то

\[
\boxed{x_2=2{,}5x_1}.
\]

Подставляем эту связь:

\[
\frac{x_1+2{,}5x_1}{2}=42.
\]

Собираем подобные слагаемые:

\[
\frac{3{,}5x_1}{2}=42.
\]

Тогда

\[
3{,}5x_1=84,
\]

\[
x_1=24.
\]

Второе число:

\[
x_2=2{,}5\cdot24=60.
\]

Ответ:

\[
\boxed{24 \text{ и } 60}.
\]

\topicsection{5. «В несколько раз больше» и «в несколько раз меньше»}

Пусть \(k>1\).

Если число \(b\) в \(k\) раз больше числа \(a\), то

\[
\boxed{b=ka}.
\]

Эквивалентная запись:

\[
\boxed{a=\frac{b}{k}}.
\]

Поэтому фраза

\[
\text{«}a\text{ в }k\text{ раз меньше }b\text{»}
\]

означает

\[
\boxed{a=\frac{b}{k}}.
\]

Например:

\[
a \text{ в }2{,}5\text{ раза меньше }b
\]

означает

\[
a=\frac{b}{2{,}5},
\]

или, что то же самое,

\[
b=2{,}5a.
\]

\subtopic{Важно различать}

\[
\text{«на }k\text{ больше»}
\quad\Longrightarrow\quad
b=a+k,
\]

\[
\text{«в }k\text{ раз больше»}
\quad\Longrightarrow\quad
b=ka.
\]

Аналогично:

\[
\text{«на }k\text{ меньше»}
\quad\Longrightarrow\quad
a=b-k,
\]

\[
\text{«в }k\text{ раз меньше»}
\quad\Longrightarrow\quad
a=\frac{b}{k}.
\]

\topicsection{6. Последовательные натуральные числа}

Если первое натуральное число обозначить через \(n\), то следующие за ним числа:

\[
\boxed{
n,\quad n+1,\quad n+2,\quad n+3,\dots
}
\]

Таким образом, три последовательных натуральных числа удобно записывать как

\[
\boxed{
n,\quad n+1,\quad n+2
}
\]

Именно выражения \(+1\) и \(+2\) показывают, что числа идут подряд.

\subtopic{Пример}

Среднее арифметическое трёх последовательных натуральных чисел равно \(21\).

Составляем модель:

\[
\frac{n+(n+1)+(n+2)}{3}=21.
\]

Упрощаем числитель:

\[
\frac{3n+3}{3}=21.
\]

Получаем

\[
n+1=21,
\]

откуда

\[
n=20.
\]

Следовательно, числа:

\[
\boxed{20,\ 21,\ 22}.
\]

\topicsection{7. Закономерность для трёх последовательных чисел}

Рассмотрим три последовательных числа:

\[
n,\qquad n+1,\qquad n+2.
\]

Их среднее арифметическое:

\[
\overline{x}
=
\frac{n+(n+1)+(n+2)}{3}.
\]

Собираем подобные слагаемые:

\[
\overline{x}
=
\frac{3n+3}{3}.
\]

Почленно делим числитель на \(3\):

\[
\overline{x}
=
\frac{3n}{3}
+
\frac{3}{3}.
\]

Получаем

\[
\boxed{
\overline{x}=n+1
}
\]

Но \(n+1\) --- среднее из трёх последовательных чисел.

Следовательно:

\[
\boxed{
\text{среднее арифметическое трёх последовательных чисел
равно среднему из них}
}
\]

Например,

\[
15,\ 16,\ 17
\]

имеют среднее арифметическое

\[
\boxed{16}.
\]

\topicsection{8. Почленное деление суммы}

Если весь числитель представляет собой сумму, то можно разделить каждое его слагаемое на знаменатель:

\[
\boxed{
\frac{a+b}{c}
=
\frac{a}{c}+\frac{b}{c},
\qquad c\ne0
}
\]

Например,

\[
\frac{3n+3}{3}
=
\frac{3n}{3}+\frac{3}{3}
=
n+1.
\]

Аналогично,

\[
\frac{a-b}{c}
=
\frac{a}{c}-\frac{b}{c},
\qquad c\ne0.
\]

\important{Делить нужно каждое слагаемое или вычитаемое числителя.}

Нельзя записывать

\[
\frac{a+b}{c}
=
a+\frac{b}{c}.
\]

\topicsection{9. Как решать задачи на среднее арифметическое}

\begin{enumerate}
    \item Определите, какие числа или группы чисел рассматриваются.

    \item Запишите основную формулу:
    \[
    \overline{x}
    =
    \frac{x_1+x_2+\dots+x_n}{n}.
    \]

    \item Переведите словесные связи на математический язык.

    \item Если нужно объединить группы, сначала восстановите суммы:
    \[
    S=n\overline{x}.
    \]

    \item Подставьте все известные величины в формулу.

    \item Сведите задачу к уравнению с одной неизвестной.

    \item Решите уравнение.

    \item Проверьте найденные значения по условию.
\end{enumerate}

\topicsection{10. Что особенно важно контролировать}

\begin{itemize}
    \item Среднее арифметическое зависит одновременно от суммы чисел и их количества.

    \item При объединении групп нужно складывать их суммы.

    \item Если группы имеют разный размер, их средние имеют разный ``вес''.

    \item Слова задачи нужно переводить в формулы до начала вычислений.

    \item Индексы \(x_1,x_2,x_3\) только нумеруют числа и сами по себе не означают, что числа последовательные.

    \item Последовательные натуральные числа записываются через прибавление единицы:
    \[
    n,\ n+1,\ n+2.
    \]

    \item Фразы ``на несколько'' и ``в несколько раз'' обозначают разные действия.

    \item Если требуется только сумма группы, искать каждое число группы отдельно обычно не нужно.

    \item При делении алгебраической суммы на число деление должно относиться ко всем слагаемым.
\end{itemize}

% ================================================================
% ТРЕНИРОВОЧНЫЙ БЛОК
% ================================================================
\newpage
\thispagestyle{fancy}

\begin{center}
    {\LARGE\bfseries\color{darkaccent}Тренировочный блок}
    
    \vspace{0.3em}
    
    {\large Среднее арифметическое и математические модели}
\end{center}

\vspace{0.5em}
{\color{accent}\rule{\linewidth}{0.8pt}}

\begin{enumerate}

    \item Первое из двух чисел равно \(8\), а их среднее арифметическое равно \(6{,}4\).
    Найдите второе число.

    \item Среднее арифметическое пяти чисел равно \(7{,}2\).
    Найдите сумму этих пяти чисел.

    \item Среднее арифметическое \(9\) чисел равно \(4{,}6\), а среднее арифметическое других \(6\) чисел равно \(7{,}1\).
    Найдите среднее арифметическое всех \(15\) чисел.

    \item Среднее арифметическое двух положительных чисел равно \(36\).
    Одно число в \(3\) раза больше другого.
    Найдите оба числа.

    \item Среднее арифметическое трёх последовательных натуральных чисел равно \(47\).
    Найдите эти числа.

    \item Первое из двух чисел равно \(12{,}8\), а их среднее арифметическое равно \(9{,}35\).
    Найдите второе число.

    \item Среднее арифметическое \(12\) чисел равно \(3{,}75\), а среднее арифметическое других \(8\) чисел равно \(6{,}25\).
    Найдите среднее арифметическое всех \(20\) чисел.

    \item Сумма трёх последовательных натуральных чисел равна \(114\).
    Найдите эти числа.

    \item Среднее арифметическое \(10\) чисел равно \(8{,}4\).
    После добавления ещё одного числа среднее арифметическое всех \(11\) чисел стало равно \(9\).
    Какое число добавили?

    \item В первой группе находится \(8\) чисел, их среднее арифметическое равно \(12{,}5\).
    Во второй группе среднее арифметическое равно \(8\).
    После объединения двух групп среднее арифметическое всех чисел стало равно \(10\).
    Сколько чисел было во второй группе?

\end{enumerate}

% ================================================================
% ОТВЕТЫ
% ================================================================
\newpage
\thispagestyle{fancy}

\begin{center}
    {\LARGE\bfseries\color{darkaccent}Ответы}
\end{center}

\vspace{0.5em}
{\color{accent}\rule{\linewidth}{0.8pt}}

\vspace{0.8em}

\renewcommand{\arraystretch}{1.4}

\begin{table}[ht]
\centering
\caption{Ответы к тренировочному блоку}
\begin{tabularx}{\linewidth}{@{}cX@{}}
\toprule
\textbf{№} & \textbf{Ответ} \\
\midrule
1 & \(4{,}8\) \\
2 & \(36\) \\
3 & \(5{,}6\) \\
4 & \(18\) и \(54\) \\
5 & \(46,\ 47,\ 48\) \\
6 & \(5{,}9\) \\
7 & \(4{,}75\) \\
8 & \(37,\ 38,\ 39\) \\
9 & \(15\) \\
10 & \(10\) чисел \\
\bottomrule
\end{tabularx}
\end{table}

\vfill

\begin{center}
{\small\color{textgray}
Лёвин Артём Александрович\\
Материал подготовлен для Нади Клименко}
\end{center}

\end{document}